\section{Controller-laget}
\label{sec:controllers}

API-et består av to controllers: \texttt{RoomsController} og \texttt{AuthController}. Begge benytter \texttt{[ApiController]}-attributtet, som gir automatisk modellvalidering og standardiserte feilsvar. Controllerne er tynne og fungerer primært som et bindeledd mellom HTTP-laget og forretningslogikken: De binder innkommende forespørsler, delegerer operasjoner til provider-laget og logger vellykkede operasjoner, uten å inneholde forretningslogikk eller \texttt{try/catch}-blokker.

\texttt{RoomsController} er beskyttet med \texttt{[Authorize]} på
klassenivå og eksponerer syv endepunkter under \texttt{/api/rooms} som dekker kravene F1.1--F1.5 og S1.2 \cite{hdo_kravspec}, samt OPTIONS-handlers for CORS-preflight. Videomotor velges via et prefiks kodet inn i ressurs-ID-en: Ved opprettelse sendes ønsket provider i request-bodyen og inkluderes i den returnerte ID-en på formen \texttt{provider:rawId}. På øvrige endepunkter parses
prefikset ut av \texttt{roomId} og brukes til å slå opp riktig
\texttt{IVideoProvider} via \texttt{IVideoProviderFactory}.
Feilhåndtering er sentralisert i \texttt{GlobalExceptionHandler}, som tilordner domenespesifikke exceptions fra provider-laget til HTTP-statuskoder og returnerer standardiserte feilsvar basert på RFC~7807 ProblemDetails-formatet. Denne ble obsoletet av RFC 9457 i 2023, men endringene er hovedsakelig redaksjonelle og påvirker ikke det vi har implementert i prosjektet.

\texttt{AuthController} har ett endepunkt, \texttt{POST /auth/token},
som henter et Bearer-token via OAuth~2.0 client credentials-flyten.
Endepunktet er annotert med \texttt{[AllowAnonymous]} ettersom det må være tilgjengelig før autentisering er etablert, og med \texttt{[EnableRateLimiting("auth")]} for å begrense antall forespørsler per minutt per IP-adresse. Selve kommunikasjonen mot Keycloak er delegert til \texttt{KeycloakService}, som sender en \texttt{POST}-forespørsel med \texttt{client\_id} og \texttt{client\_secret} til Keycloaks \texttt{/protocol/openid-connect/token}-endepunkt og returnerer access-tokenet. Tjenesten er stateless og uten caching; ansvaret for gjenbruk av token ligger hos kalleren.